\section*{Bab \ref{ch:g9:unique-individuals} --- Individu yang Unik}

\begin{solution}{exo:g9:unique-individuals:1}
Pengocokan pada pembagian duanya (beserta pertukaran
ruasnya), lotere \omterm{def:g5:flower-to-fruit:steps}{pembuahannya}, dan tetesan \omterm{prop:g9:genes-and-dna:explains}{mutasinya} dari
angkatan ke angkatan.
\end{solution}

\begin{solution}{exo:g9:unique-individuals:2}
Ia memotong tumpukan itu di bawah \omterm{def:g9:chromosomes:chromosomes}{kromosom} yang utuh: mitranya
bertukar ruas sebelum memisahkan diri, sehingga bahkan satu
\omterm{def:g9:chromosomes:chromosomes}{kromosom} pun meneruskan sebuah mozaik yang baru --- bilangan
$2^{23}$ hanyalah lantainya.
\end{solution}

\begin{solution}{exo:g9:unique-individuals:3}
Dua manusia mana pun sepakat pada lebih dari $99$ persen
teks \omterm{def:g9:genes-and-dna:dna}{DNA-nya}, dan ketidaksepakatan yang tersebar itu tetap
membuat tiap individu unik secara ketat --- \omterm{prop:g6:classification-kinship:kinship}{kekerabatan} dan
perbedaan di dalam satu molekul.
\end{solution}

\begin{solution}{exo:g9:unique-individuals:4}
Kehidupannya memencar sesudah pembagian bersamanya: \omterm{def:g9:heredity-and-traits:traits}{sifat}
perolehan, \omterm{def:g8:nervous-communication:synapse}{sinapsis} yang terlatih, dan \omterm{prop:g9:genes-and-dna:explains}{mutasi} lanjut tubuh
masing-masing menuliskan catatan yang dapat dibedakan di atas
permulaan yang serupa.
\end{solution}

\begin{solution}{exo:g9:unique-individuals:5}
\omterm{def:g5:first-look-at-cells:cell}{Sel} siapa sebuah jejaknya (pengenalannya), dan siapa orang
tua seorang anak (tiap kedudukan yang berubah-ubah harus
terbagikan dari kedua tumpukan orang tuanya).
\end{solution}

\begin{solution}{exo:g9:unique-individuals:6}
Karena sistem kekebalannya membaca penanda-diri yang halus
butirannya, dan makin dekat \omterm{prop:g6:classification-kinship:kinship}{kekerabatannya}, makin dekat
pembagian penandanya --- tumpukan keluarganya paling banyak
berbagi.
\end{solution}

\begin{solution}{exo:g9:unique-individuals:7}
Tiap orang tua menawarkan jutaan setengah tumpukan yang
mungkin; seorang anak adalah satu pemasangan satu pembagian
dari ibunya dengan satu pembagian dari ayahnya, dipilih oleh
\omterm{def:g8:sexual-reproduction:gametes}{gamet} mana yang menemui \omterm{def:g8:sexual-reproduction:gametes}{gamet} mana --- kedua ruangnya
mengalikan diri, dan hasil kalinya begitu melampaui seluruh
manusia yang pernah lahir sehingga sebuah pembagian yang
berulang boleh dibilang tidak pernah terjadi dan tidak akan
pernah terjadi.
\end{solution}

\begin{solution}{exo:g9:unique-individuals:8}
Sebuah cadangan: klonanya adalah satu tumpukan yang diulang
--- satu susunan tubuh, satu kelemahan bersama --- sementara
tumpukan unik sebuah \omterm{def:g8:reproduction-and-environment:population}{populasi} yang tumbuh dari \omterm{def:g2:from-seed-to-plant:seed}{biji} akan
menyebarkan ketahanan di antara mereka, dan sebagian
pohonnya bertahan di tempat hamanya menemukan sisanya sesuai
seleranya.
\end{solution}

\begin{solution}{exo:g9:unique-individuals:9}
Tingkatnya: klaimnya melompat dari sebuah \omterm{def:g9:genes-and-dna:gene}{gen} ke sebuah
hasil kehidupan setingkat keluarga --- perbesarannya
dilewati. Tenunannya: keberhasilan sekolah adalah rentang
turunan yang dijalani di sekolah, rumah dan bahasanya ---
satu \omterm{def:g9:genes-and-dna:gene}{gen} sama dengan tidak ada takdir. Keragamannya: ia
menilai sebuah sebaran terhadap sebuah ``normal'' yang
hantu. Kembarnya: tumpukan yang serupa lazim menghasilkan
kisah sekolah yang berbeda. Tajuk beritanya gagal pada
keempat pemeriksaannya.
\end{solution}

\begin{solution}{exo:g9:unique-individuals:10}
Karena spesiesnya \emph{adalah} keragamannya: cadangannya
menyebarkan \omterm{def:g9:genes-and-dna:gene}{alel} ke seluruh individu yang unik, dan tidak
ada tumpukan yang menjadi baku, yang darinya yang lain
menyimpang --- ``normal'' menamai sebuah rata-rata yang tak
diangkut siapa pun, bukan sebuah teks yang seharusnya
dicocoki siapa pun.
\end{solution}

\begin{solution}{exo:g9:unique-individuals:11}
Pada tiap kedudukan yang berubah-ubah, kedua ejaan anaknya
harus datang satu dari pasangan ibunya, satu dari pasangan
ayahnya. Ujinya memeriksa kedudukan demi kedudukan: tiap
ejaan anak bermata birunya dapat dibagikan dari kedua
tumpukan itu --- lompatannya menjelaskan \omterm{def:g1:the-five-senses:sense}{matanya}, kedudukannya
mengesahkan keayahannya.
\end{solution}

\begin{solution}{exo:g9:unique-individuals:12}
Misalnya: satu pengocokan membagikan kepada tiap manusia
sebuah pilihan yang tak pernah berulang dari tumpukan orang
tuanya, dan \omterm{prop:g9:genes-and-dna:explains}{mutasi} menggarami tiap angkatan dengan ejaan
yang sama sekali belum pernah terlihat --- karena itu
semuanya berbeda. Namun setiap tumpukan tertulis dengan
empat huruf yang sama, dibaca oleh sandi yang sama, dan
sepakat dengan tumpukan manusia mana pun yang lain pada
lebih dari $99$ huruf tiap seratusnya --- karena itu semuanya
sekerabat. Mekanismenya bukanlah lawan melainkan satu
sistem: teks yang sama, dibagikan ulang tanpa henti ---
\omterm{prop:g6:classification-kinship:kinship}{kekerabatannya} memasok tumpukannya, pengocokannya menjamin
tidak ada tangan yang pernah berulang.
\end{solution}

\begin{solution}{pb:g9:unique-individuals:1}
\textbf{1.} Karena \omterm{def:g9:genes-and-dna:dna}{DNA} tinggal di dalam \omterm{def:g6:cell-unit-of-life:parts}{inti selnya}, dan
batang rambutnya adalah keratin yang mati --- \omterm{def:g5:first-look-at-cells:cell}{sel} akarnyalah
yang mengangkut \omterm{def:g6:cell-unit-of-life:parts}{inti sel} yang dibutuhkan pembacaan lengkapnya.

\textbf{2.} Kedudukan yang berubah-ubah adalah tempat
individunya berbeda --- sebaran yang mengenali. Kedudukan
bersamanya hanya akan menegakkan bahwa rambutnya milik
manusia: kebenaran setingkat \omterm{prop:g6:classification-kinship:kinship}{kekerabatan}, tak berguna untuk
menyebutkan satu orang.

\textbf{3.} Bahwa \omterm{def:g5:first-look-at-cells:cell}{sel} rambutnya dan \omterm{def:g5:first-look-at-cells:cell}{sel} tersangkanya berbagi
satu pembagian --- satu individu, dari semua manusia yang
pernah ada, dengan satu-satunya pengecualian seorang kembar
identik, yang berbagi pembagiannya sejak asalnya.

\textbf{4.} Menyingkirkan atau menyertakan kembarnya dengan
bukti yang lain --- keberadaannya, saksinya, ``bagaimana ia
sampai ke sana'' pada perkaranya: profilnya sungguh tak dapat
memisahkan keduanya, sehingga pengadilannya yang harus.

\textbf{5.} Seorang kerabat berbagi, lewat keturunan,
sebagian besar ejaan yang berubah-ubahnya --- lebih banyak
daripada seorang asing, lebih sedikit daripada sebuah
kecocokan: tumpang tindih sebagian yang berjenjang menurut
\omterm{prop:g6:classification-kinship:kinship}{kekerabatan} adalah pembagiannya, dijalankan mundur menaiki
pohonnya.

\textbf{6.} Tingkatnya: sebuah profil adalah data setingkat
tumpukan --- kedudukan berubah-ubah yang dibaca justru
dipilih karena tidak mengatakan apa pun tentang \omterm{def:g9:heredity-and-traits:traits}{sifatnya}.
Tenunannya: perangai dan kelakuan adalah rentang yang
dijalani di dalam kehidupan. Keragamannya: ``\omterm{def:g9:genes-and-dna:gene}{gen} penjahat''
menilai orang terhadap sebuah normal yang hantu. Kembarnya:
tumpukan yang sama, kehidupan yang berbeda. Profilnya
menyebutkan \omterm{def:g5:first-look-at-cells:cell}{sel}, bukan watak.

\textbf{7.} Jumlah kedudukan berubah-ubah yang saling bebas:
tiap kedudukan yang ditambahkan membagi lebih jauh \omterm{def:g8:reproduction-and-environment:population}{populasi}
yang cocok --- dengan cukup banyak kedudukan, ``sekali di
kota ini'' menjadi sekali di seluruh umat manusia.
Laboratoriumnya membaca sebanyak yang dituntut bakunya.

\textbf{8.} \omterm{def:g1:living-or-not:living}{Biologinya} menegaskan \omterm{def:g5:first-look-at-cells:cell}{sel} siapa yang menumbuhkan
rambutnya. Perkaranya masih harus menegakkan penempatannya
--- kapan, bagaimana, dengan polos atau tidak, rambutnya
tiba: kecocokan bukanlah cerita, dan pengadilan menghukum
atas cerita yang terbukti.

\textbf{9.} Pembagian pembagian-dua-dan-pembuahannya membuat
tiap tumpukan tak terulang, dan perbedaan ejaan \omterm{def:g9:genes-and-dna:dna}{DNA} membuat
pembagiannya terbaca --- \omterm{def:g9:chromosomes:chromosomes}{kromosomnya} mengocoknya, molekulnya
mencatatnya.

\textbf{10.} Pada kemantapan kimiawi \omterm{def:g9:genes-and-dna:dna}{DNA}: teksnya tidak
memudar dan tidak menulis ulang dirinya di dalam sebuah laci
--- untai ganda yang sama, utuh sesudah tiga puluh tahun,
terbaca sebagaimana ia dibagikan.

\textbf{11.} Yang melayani: pengenalan yang membersihkan
nama orang yang tak bersalah dan memecahkan soal keayahan;
pencocokan cangkok yang menemukan pembagian yang terdekat.
Penyalahgunaan yang diperingatkan: membaca keragaman
spesiesnya sebagai peringkat --- perbedaan yang dipaksa
menjadi tingkatan, \omterm{def:g1:living-or-not:living}{biologi} buruk di balik ketidakadilan
zaman dahulu.

\textbf{12.} Tidak ada dua \omterm{def:g5:flower-to-fruit:steps}{pembuahan} yang pernah membagikan
tumpukan yang sama: ruang setengah tumpukan orang tuanya
dikalikan lotere \omterm{def:g8:sexual-reproduction:gametes}{gametnya} melampaui seluruh manusia yang
pernah dikandung --- sehingga sebuah kecocokan profil penuh
menyebutkan satu individu, dan pengadilannya boleh
memperlakukan rambutnya sebagai tertandatangani.
\end{solution}
